\begin{table}[htbp]
\centering
\caption{Apoio à produção, acesso a programas e a mercado, por subpopulação}
\label{tab:apoio_sub}
\resizebox{\textwidth}{!}{%
\begin{tabular}{l r r r r r }
\hline
\textbf{Indicador (\%)} & Agr. Familiar & Agr. Empresarial & Assentamento & Moradia e Lazer & \textbf{Total} \\ \hline
Incentivo à produção (já recebeu) & 9,1 & 6,0 & 8,4 & 3,2 & 5,5 \\
\quad Pronaf (entre os beneficiários) & 24,7 & 39,5 & 33,7 & 36,9 & 32,3 \\
\quad ATER (entre os beneficiários) & 28,5 & 26,1 & 13,3 & 13,6 & 22,9 \\
\quad Fomento rural (entre os beneficiários) & 5,6 & 1,1 & 18,0 & 5,0 & 5,0 \\
\quad FCO (entre os beneficiários) & 3,4 & 1,6 & 0,0 & 3,2 & 2,7 \\
Outorga de uso da água & 44,4 & 41,3 & 48,8 & 34,8 & 39,6 \\
Benefício de programa social & 9,5 & 8,6 & 11,6 & 11,3 & 10,3 \\
Pediu auxílio governamental (via internet) & 9,3 & 11,8 & 7,2 & 6,6 & 8,4 \\
Possui CNPJ (ocupados) & 7,5 & 8,6 & 5,6 & 7,6 & 7,7 \\
Transações financeiras pela internet (PIX) & 48,5 & 52,4 & 48,8 & 50,8 & 50,5 \\
Curso profissional agrícola em andamento & 4,2 & 2,9 & 1,6 & 3,2 & 3,3 \\
\hline
Base: domicílios com resposta válida & 1.355 & 1.168 & 618 & 1.785 & 4.926 \\
Base: domicílios beneficiários & 130 & 71 & 63 & 54 & 318 \\
\hline
\end{tabular}}
\par\footnotesize Fonte: elaboração própria a partir da PDAD Rural 2022. Proporções ponderadas pelo peso corrigido, entre as respostas válidas (``sim'' ou ``não''); respostas ``não sabe'' e ``não se aplica'' excluídas. Os programas são categorias excludentes entre os domicílios que receberam incentivo (as demais respostas são ``outro'' e ``não sabe''). $^{a}$ Base de menos de 30 domicílios beneficiários: leitura apenas indicativa.
\end{table}
